\subsection{Restriction of differentiable sections of a vector bundle}

Let r be a positive integer, X a differentiable manifold of class \(C^{r}\) locally of finite dimension, and E a vector bundle (real or complex),

with base X and of class \(C^{r}\). For every open set U of X, denote by \(\mathcal{S}^{r}(\mathrm{U};\mathrm{E})\) the vector space (denoted \(\mathcal{S}_{\mathrm{E}}^{r}(\mathrm{U})\) in VAR, R1, 7.4, p. 74) of \(C^{r}\) sections of E over U, equipped with the topology of \(C^{r}\) compact convergence.

Lemma 4. --- Let \(\mathcal{U}\) be an open cover of X. The map \(u: f \mapsto (f|U)_{U \in \mathcal{U}}\) from \(\mathcal{S}^{r}(X;E)\) into \(\prod_{U \in \mathcal{U}} \mathcal{S}^{r}(U;E)\) is linear, injective, continuous, and strict. Its image is closed.

The map \(u\) is linear, injective, and continuous. By virtue of TG, IX, p. 43, prop. 1 and p. 48, cor. 1, every compact subset of X has a finite cover \((\mathrm{C}_i)_{i\in \mathrm{I}}\) such that, for each \(i\in \mathrm{I}\) , the set \(\mathrm{C}_i\) is a compact subset of one of the open sets of the covering \(\mathcal{U}\) . It follows that the linear map \(u\) is strict. Finally, its image is closed, since it consists of the families \((f_{\mathrm{U}})_{\mathrm{U}\in \mathcal{U}}\) such that \(f_{\mathrm{U}}\) and \(f_{\mathrm{V}}\) coincide in \(\mathrm{U}\cap \mathrm{V}\) for all U and V in \(\mathcal{U}\) .

PROPOSITION 9. --- The space \(\mathcal{S}^r(\mathrm{X};\mathrm{E})\) is complete.

Suppose first that there exists an integer \(n \geqslant 0\) and a Banach space F such that X is an open subset of \(R^{n}\) and E is the trivial vector bundle \(X \times F\) with base X and fiber F. In this case, the topological vector space \(\mathcal{S}^{r}(X;E)\) is isomorphic to \(\mathcal{C}^{r}(X;F)\) , and the result follows from lemma 3 of III, p. 34.

In the general case, let \(\mathcal{U}\) be an open covering of X by coordinate domains such that for every \(\mathrm{U} \in \mathcal{U}\) , the restriction of E to U is trivializable. By the lemma above, the space \(\mathcal{S}^r(\mathrm{X};\mathrm{E})\) is isomorphic to the image of the linear map \(f \mapsto (f|\mathrm{U})_{\mathrm{U} \in \mathcal{U}}\) , which is closed in the product of the spaces \(\mathcal{S}^r(\mathrm{U};\mathrm{E})\) for U in \(\mathcal{U}\) . By the preceding case, each of the spaces \(\mathcal{S}^r(\mathrm{U};\mathrm{E})\) is complete, hence their product is complete (TG, II, p. 17, prop. 10). Consequently, \(\mathcal{S}^r(\mathrm{X};\mathrm{E})\) is complete.

Remark. --- Let \(\mathcal{U}\) be an open covering of X by coordinate domains \(c_{\mathrm{U}} = (\mathrm{U}, \varphi_{\mathrm{U}}, \mathbf{R}^{n_{\mathrm{U}}})\) such that for every \(U \in \mathcal{U}\) , the restriction of E to U is trivializable of type \(F_{U}\) , where \(F_{U}\) is a Banach space. For every \(U \in \mathcal{U}\) , one identifies a section f of E over U with a map \(f_{U}\) from \(\varphi_{\mathrm{U}}(\mathrm{U})\) into \(F_{U}\) . It follows from the proof of prop. 9 that the topology of the space \(\mathcal{S}^{r}(\mathrm{X};\mathrm{E})\) is defined by the family of seminorms \(p_{U,C,\alpha}\) such that

\[
p _ {\mathrm{U,C}, \alpha} (f) = \sup _ {x \in \mathrm{C}} \| (\partial^ {\alpha} f _ {\mathrm{U}}) (x) \|,
\]

where U ranges over \(\mathcal{U}\) , C ranges over the set of compact subsets of \(\varphi_{\mathrm{U}}(\mathrm{U})\) and \(\alpha \in \mathbf{N}^{n_{\mathrm{U}}}\) ranges over the set of multi-indices such that \(|\alpha| \leqslant r\) .

PROPOSITION 10. --- Suppose that the vector bundle E is locally of finite rank. Let s be an integer such that \(0 \leqslant s < r\) and Y a relatively compact open subset of X. The linear map \(f \mapsto f|Y\) from \(\mathcal{S}^{r}(X;E)\) into \(\mathcal{S}^{s}(Y;E)\) is compact.

Suppose first that there exists an integer \(n \geqslant 0\) and a finite-dimensional Banach space F such that X is an open subset of \(R^{n}\) and E is the trivial vector bundle \(X \times F\) with base X and fiber F. In this case, the topological vector spaces \(\mathcal{S}^{r}(X;E)\) and \(\mathcal{S}^{s}(Y;E)\) are isomorphic to \(\mathcal{C}^{r}(X;F)\) and \(\mathcal{C}^{s}(Y;F)\) , respectively, and prop. 10 is a consequence of prop. 8.

Let us pass to the general case. Let C be a compact subset of X containing Y. For every point x of C, choose an open neighborhood \(U(x)\) of x in X which is a coordinate domain over which the vector bundle E is trivializable. Choose moreover a relatively compact open neighborhood \(V(x)\) of x in \(U(x)\) . Since the set C is compact, it is covered by a finite family \((\mathrm{V}(x_{1}),\ldots,\mathrm{V}(x_{m}))\) of such open sets. For every i, set \(U_{i}=U(x_{i})\) and \(Y_{i}=V(x_{i})\cap Y\) . Then \(Y_{i}\) is a relatively compact open subset of \(U_{i}\) and we have \(Y=Y_{1}\cup\cdots\cup Y_{m}\) . The linear maps \(f\mapsto f|U_{i}\) from \(\mathcal{S}^{r}(X;E)\) into \(\mathcal{S}^{r}(U_{i};E)\) are continuous and the linear maps \(g\mapsto g|Y_{i}\) from \(\mathcal{S}^{r}(U_{i};E)\) into \(\mathcal{S}^{s}(Y_{i};E)\) are compact by the first part of the proof. The linear maps \(f\mapsto f|Y_{i}\) from \(\mathcal{S}^{r}(X;E)\) into \(\mathcal{S}^{s}(Y_{i};E)\) are therefore compact (prop. 3 of III, p. 5). Taking into account Lemma 4, applied to the covering of Y by the open sets \(Y_{i}\) , and Remarks 5 and 6 of III, p. 3, the linear map \(f\mapsto f|Y\) from \(\mathcal{S}^{r}(X;E)\) into \(\mathcal{S}^{s}(Y;E)\) is compact, which completes the proof.

